% Part: first-order-logic
% Chapter: syntax-and-semantics
% Section: unique-readability

\documentclass[../../../include/open-logic-section]{subfiles}

\begin{document}

\olfileid[hi]{fol}{syn}{unq}

\olsection{अद्वितीय पठनीयता}

\begin{explain}
!!{formula}s की हमारी परिभाषा यह सुनिश्चित करती है कि प्रत्येक !!{formula}
का \emph{अद्वितीय पाठ} हो; अर्थात् !!{formula}s के हमारे निर्माण-नियमों के
अनुसार उसे बनाने का मूलतः केवल एक ढंग हो और उसकी ``व्याख्या'' करने का भी
केवल एक ढंग हो। ऐसा न होता तो कुछ !!{formula}s संदिग्ध होते, अर्थात् उनका
एक से अधिक पाठ या अर्थ-निर्णय होता---जिससे हम स्पष्टतः बचना चाहते हैं। इससे
भी अधिक महत्त्वपूर्ण यह है कि इस गुण के बिना आगे दी जाने वाली हमारी अधिकांश
परिभाषाएँ और प्रमाण काम नहीं करेंगे।

इसे स्पष्ट करने का शायद सर्वोत्तम ढंग यह देखना है कि यदि हमने !!{formula}s
के ऐसे दोषपूर्ण निर्माण-नियम दिए होते जो अद्वितीय पठनीयता सुनिश्चित न करते,
तो क्या होता। उदाहरणतः हम संयोजकों के निर्माण-नियमों में कोष्ठक भूल सकते थे
और यह अनुमति दे सकते थे:
\begin{quote}
यदि $!A$ और $!B$ !!{formula}s हैं, तो $!A \lif !B$ भी !!{formula} है।
\end{quote}
किसी परमाण्विक सूत्र $!D$ से आरंभ करके इससे हम $!D
\lif !D$ बना सकते। इससे और $!D$ से हमें $!D \lif !D
\lif !D$ मिलता। किंतु इसे बनाने के दो ढंग हैं:
\begin{enumerate}
\item हम $!D$ को $!A$ और $!D \lif !D$ को $!B$ लेते हैं।
\item हम $!A$ को $!D \lif !D$ और $!B$ को~$!D$ लेते हैं।
\end{enumerate}
तदनुसार !!{formula}~$!D \lif !D \lif !D$ को ``पढ़ने'' के भी दो ढंग हैं।
यह $!B
\lif !C$ के रूप में है, जहाँ $!B$, $!D$ है और $!C$, $!D \lif !D$
है; किंतु यह $!B \lif !C$ के रूप में \emph{भी है}, जहाँ $!B$,
$!D \lif !D$ है और $!C$,~$!D$ है।

यदि ऐसा हो, तो हमारी परिभाषाएँ सदा काम नहीं करेंगी। उदाहरणतः किसी सूत्र का
!!{main operator} परिभाषित करते समय हम कहते हैं कि $!B \lif !C$ रूप के सूत्र
में~$\lif$ की दर्शाई गई उपस्थिति ही !!{main operator} है। किंतु यदि सूत्र $!D
\lif !D \lif !D$ को ऊपर बताए दो भिन्न ढंगों से $!B \lif !C$ के साथ मिलाया
जा सके, तो एक स्थिति में $\lif$ की पहली उपस्थिति !!{main operator} मिलती है
और दूसरी स्थिति में दूसरी उपस्थिति। हमारा आशय तो यह है कि !!{main operator},
!!{formula} का कोई \emph{फलन} हो; अर्थात् प्रत्येक !!{formula} में
!!{main
  operator} की ठीक एक उपस्थिति होनी चाहिए।
\end{explain}

\begin{lem}
किसी !!a{formula}~$!A$ में बाएँ और दाएँ कोष्ठकों की संख्या समान होती है।
\end{lem}

\begin{proof}
हम इसका प्रमाण $!A$ के निर्माण पर आगमन द्वारा देते हैं। इसके लिए दो बातें
चाहिए: (क) पहले सिद्ध करना होगा कि सभी परमाण्विक सूत्रों में यह गुण है
(आगमन-आधार)। (ख) फिर सिद्ध करना होगा कि दिए गए सूत्रों से नया सूत्र बनाने पर,
यदि पुराने सूत्रों में यह गुण है तो नए सूत्र में भी है।

मान लें कि $l(!A)$ बाएँ कोष्ठकों की संख्या और $r(!A)$ दाएँ कोष्ठकों की
संख्या है जो~$!A$ में हैं; इसी प्रकार $l(t)$ और $r(t)$ किसी पद~$t$ में
क्रमशः बाएँ और दाएँ कोष्ठकों की संख्याएँ हैं।

\begin{prob}
    सिद्ध कीजिए कि किसी भी पद~$t$ के लिए $l(t) = r(t)$।
\end{prob}

\begin{enumerate}
\tagitem{prvFalse}{\indcase{!A}{\lfalse}{$\indfrm$ में $0$ बाएँ और $0$
    दाएँ कोष्ठक हैं।}}{}

\tagitem{prvTrue}{\indcase{!A}{\ltrue}{$\indfrm$ में $0$ बाएँ और $0$
    दाएँ कोष्ठक हैं।}}{}

\item \indcase{!A}{\Atom{R}{t_1,\dots,t_n}}{$l(\indfrm) = 1 + l(t_1) +
  \dots + l(t_n) = 1 + r(t_1) + \dots + r(t_n) = r(\indfrm)$। यहाँ हम
  इस तथ्य का उपयोग करते हैं कि $l(t) = r(t)$ किसी भी पद~$t$ के लिए; इसे
  अभ्यास के लिए छोड़ा गया है।}

\item \indcase{!A}{\eq[t_1][t_2]}{$l(\indfrm) = l(t_1) + l(t_2) =
  r(t_1) + r(t_2) = r(\indfrm)$।}

\tagitem{prvNot}{\indcase{!A}{\lnot !B}{आगमनात्मक परिकल्पना से
    $l(!B) = r(!B)$। अतः $l(\indfrm) = l(!B) = r(!B) =
    r(\indfrm)$।}}{}

\item \indcase{!A}{(!B \ast !C)}{आगमनात्मक परिकल्पना से $l(!B) =
  r(!B)$ और $l(!C) = r(!C)$। अतः $l(\indfrm) = 1 + l(!B) + l(!C) =
  1 + r(!B) + r(!C) = r(\indfrm)$।}

\tagitem{prvAll}{\indcase{!A}{\lforall[x][!B]}{आगमनात्मक परिकल्पना से
    $l(!B) = r(!B)$। अतः $l(\indfrm) = l(!B) = r(!B) =
    r(\indfrm)$।}}{}

% Print case for \lexists if prvEx and not prvAll
\tagitem{prvAll,notprvEx}{}{\indcase{!A}{\lexists[x][!B]}{आगमनात्मक
    परिकल्पना से $l(!B) = r(!B)$। अतः $l(\indfrm) = l(!B) = r(!B) =
    r(\indfrm)$।}}

% Just say similarly if prvEx and prvAll
\tagitem{notprvAll,notprvEx}{}{\indcase{!A}{\lexists[x][!B]}{इसी प्रकार।}}
\end{enumerate}
\end{proof}

\begin{defn}[उचित आरंभिक खंड]
प्रतीक-शृंखला $!B$ किसी प्रतीक-शृंखला~$!A$ का \emph{उचित आरंभिक खंड} है यदि
$!B$ और किसी अरिक्त प्रतीक-शृंखला को जोड़ने पर~$!A$ प्राप्त हो।
\end{defn}

\begin{lem}\ollabel{lem:no-prefix}
यदि $!A$ !!a{formula} है और $!B$, $!A$ का उचित आरंभिक खंड है, तो $!B$
!!a{formula} नहीं है।
\end{lem}

\begin{proof}
अभ्यास।
\end{proof}

\begin{prob}
\olref[fol][syn][unq]{lem:no-prefix} सिद्ध कीजिए।
\end{prob}

\begin{prop}
\ollabel{prop:unique-atomic}
यदि $!A$ कोई परमाण्विक !!{formula} है, तो वह निम्न शर्तों में से ठीक एक को
संतुष्ट करता है।
\begin{enumerate}
\tagitem{prvFalse}{$!A \ident \lfalse$.}{}
\tagitem{prvTrue}{$!A \ident \ltrue$.}{}
\item $!A \ident \Atom{R}{t_1,\dots,t_n}$, जहाँ $R$ कोई $n$-स्थानी
  !!{predicate} है, $t_1$, \dots, $t_n$ पद हैं, और $R$, $t_1$, \dots,
  $t_n$ में से प्रत्येक अद्वितीय रूप से निर्धारित है।
\item $!A \ident \eq[t_1][t_2]$, जहाँ $t_1$ और $t_2$ अद्वितीय रूप से
  निर्धारित पद हैं।
\end{enumerate}
\end{prop}

\begin{proof}
अभ्यास।
\end{proof}

\begin{prob}
\olref[fol][syn][unq]{prop:unique-atomic} सिद्ध कीजिए। (संकेत: पदों के लिए
\olref[fol][syn][unq]{lem:no-prefix} का एक रूप बनाकर सिद्ध कीजिए।)
\end{prob}

\begin{prop}[अद्वितीय पठनीयता]
प्रत्येक !!{formula} निम्न शर्तों में से ठीक एक को संतुष्ट करता है।
\begin{enumerate}
\item $!A$ परमाण्विक है।

\tagitem{prvNot}{$!A$, $\lnot !B$ के रूप में है।}{}

\tagitem{prvAnd}{$!A$, $(!B \land !C)$ के रूप में है।}{}

\tagitem{prvOr}{$!A$, $(!B \lor !C)$ के रूप में है।}{}

\tagitem{prvIf}{$!A$, $(!B \lif !C)$ के रूप में है।}{}

\tagitem{prvIff}{$!A$, $(!B \liff !C)$ के रूप में है।}{}

\tagitem{prvAll}{$!A$, $\lforall[x][!B]$ के रूप में है।}{}

\tagitem{prvEx}{$!A$, $\lexists[x][!B]$ के रूप में है।}{}
\end{enumerate}
इसके अतिरिक्त प्रत्येक स्थिति में $!B$, अथवा $!B$ और $!C$, अद्वितीय रूप से
निर्धारित होते हैं। इसका अर्थ है कि, उदाहरणतः, ऐसे भिन्न युग्म $!B$, $!C$
और $!B'$, $!C'$ नहीं हैं जिनसे $!A$ एक साथ इन दोनों रूपों में हो:
\iftag{prvIf}{$(!B \lif !C)$ और $(!B' \lif !C')$।}{
    \iftag{prvOr}{$(!B \lor !C)$ और $(!B' \lor !C')$।}{
      \iftag{prvAnd}{$(!B \land !C)$ और $(!B' \land !C')$।}{}}}
\end{prop}

\begin{proof}
निर्माण-नियमों के अनुसार यदि कोई !!a{formula} परमाण्विक नहीं है, तो उसे
आरंभिक कोष्ठक~(, \iftag{prvNot}{$\lnot$,}{} या किसी परिमाणक से आरंभ होना
चाहिए। दूसरी ओर, निम्न प्रतीकों में से किसी से आरंभ होने वाला प्रत्येक
!!{formula} परमाण्विक होना चाहिए: !!a{predicate}, !!a{function},
!!a{constant}\iftag{prvFalse}{, $\lfalse$}{}\iftag{prvTrue}{, $\ltrue$}{}।

अतः हमें वास्तव में केवल यह दिखाना है कि यदि $!A$, $(!B \ast
!C)$ के रूप
में और $(!B' \mathbin{\ast'} !C')$ के रूप में भी है, तो $!B \ident
!B'$,
$!C \ident !C'$, और $\ast = {\ast'}$।

मान लें कि $!A \ident (!B \ast !C)$ और $!A \ident (!B'
\mathbin{\ast'} !C')$ दोनों हैं। तब या तो $!B \ident !B'$ है या नहीं। यदि
है, तो स्पष्टतः $\ast = {\ast'}$ और $!C \ident !C'$ हैं, क्योंकि तब वे
$!A$ की ऐसी उपशृंखलाएँ हैं जो एक ही स्थान पर आरंभ होती और समान लंबाई की हैं।
दूसरी स्थिति $!B \not\ident !B'$ है। चूँकि $!B$ और $!B'$ दोनों $!A$ की
एक ही स्थान पर आरंभ होने वाली उपशृंखलाएँ हैं, उनमें से एक दूसरे का उचित
आरंभिक खंड होना चाहिए। किंतु \olref{lem:no-prefix} के अनुसार यह असंभव है।
\end{proof}


\end{document}
